% ======================================================================
% Appendix A -- Proofs
% ======================================================================
\section{Proofs}
\label{app:proofs}

\subsection{Proof of Proposition T1 (Non-Identifiability of Substitutes)}
\label{app:proof-t1}

% TODO(researcher): finalize statement + proof.
\todoresearcher{Finalize the exact statement of Proposition T1
(masked-coalition value functions cannot distinguish certain
neighborhood substitutes without retraining) and write the proof.
Sketch: construct two parameterizations of the brief run realizing the
same
function with credit routed through either substitute; show
$v_{\mathrm{prune}}$ is identical on all coalitions for one
parameterization but not under retraining.}

\begin{proof}[Proof sketch (placeholder)]
TBD. % TODO(researcher)
\end{proof}

\subsection{Statement and Proof Sketch of Proposition T3 (Rank Floor)}
\label{app:proof-t3}

% TODO(researcher): finalize statement + proof.
\todoresearcher{Finalize the exact statement of Proposition T3 (the
rank floor) and write the full proof.}

\begin{theorybox}
\begin{propositionTrank}[Rank floor; statement to be finalized]
% (no \label: unnumbered T-statement; refer to it as "Proposition~T3")
Let $\mathcal{C}$ be a complex with $R$ ranks and $S$ a coalition of
neighborhoods. If the rank pairs connected by the neighborhoods in $S$
do not span the rank hierarchy, then every cell of each unspanned rank
has an empty composed receptive field: no input signal can reach it
through any sequence of message-passing layers, and each such cell is a
null player of the masked game. Spanning $R$ ranks requires on the
order of $R$ neighborhoods; hence $k \ge R$ is a structural floor on
the size of any coalition that uses the whole complex---$k = 3$ for the
2-complexes (nodes, edges, faces) of molecular benchmarks. \emph{Exact
statement TBD by researcher.}
\end{propositionTrank}
\end{theorybox}

\begin{proof}[Proof sketch]
Each neighborhood is an operator between one pair of ranks
(\cref{sec:prelim-gccn}), so a coalition $S$ induces a
rank-connectivity multigraph $G_S$ on the ranks of the complex, with
one edge per neighborhood in $S$. Composed receptive fields grow only
along walks of $G_S$: after any number of message-passing layers, a
cell of rank $r$ can have received signal only from cells whose ranks
lie in the connected component of $r$ in $G_S$. If a rank is
disconnected from every rank carrying input signal, all of its cells
have empty composed receptive fields; masking such a cell changes no
coalition value, so it is a null player and the null-player axiom
assigns it exactly zero. Connecting all ranks of an $R$-rank complex
requires a connected spanning subgraph of $G_S$, hence on the order of
$R$ neighborhood edges; accounting for the directionality of routes
and for which ranks carry input signal versus readout yields the floor
$k \ge R$ stated above and discussed in \cref{app:selection-rankfloor}
($k \ge 3$ for the 2-complexes of the molecular benchmarks).
% TODO(researcher): tighten the directionality/readout counting into
% the final constant.
\end{proof}

\subsection{Proof of Corollary T2 (Hypergraph Message Passing)}
\label{app:proof-t2}

% TODO(researcher): finalize statement + proof.
\todoresearcher{Finalize the exact statement of Corollary T2 (the
framework extends to hypergraph message passing; carried by theory only,
the hypergraph experiment was cut) and write the proof. Sketch:
hypergraph incidence defines the neighborhood set; players, masking, and
both value functions are unchanged; axioms follow from
\cref{eq:shapley,eq:interaction} verbatim.}

\begin{proof}[Proof sketch (placeholder)]
TBD. % TODO(researcher)
\end{proof}
